\chapter{JUSTIFICATIVA}
\label{cap:justificativa}

O mercado global de jogos digitais mantém uma participação econômica expressiva no setor de entretenimento. Segundo a \citeonline{newzoo2025global}, a receita mundial do segmento foi estimada em US\$ 188,8 bilhões para 2025, representando um crescimento de 3,4\% em relação ao ano anterior. No cenário brasileiro, o mercado de jogos digitais foi estimado em cerca de US\$ 2,3 bilhões em receitas no ano de 2021, posicionando o país entre os principais mercados do setor \cite{fortim2022pesquisa}. Além de sua relevância econômica, o setor contribui para a geração de empregos qualificados e para o desenvolvimento de competências técnicas e criativas.

No entanto, a gestão do desenvolvimento de jogos enfrenta desafios recorrentes relacionados à comunicação, à baixa formalização de processos e à ausência de uma linguagem de design padronizada, fatores que dificultam a coordenação das equipes, a identificação precoce de problemas e a tomada de decisões durante o desenvolvimento \cite{giri2022exploratory}. Nesse contexto, a plataforma proposta busca contribuir para a organização de \textit{feedbacks} dispersos, permitindo que essas informações sejam consultadas e analisadas de maneira mais estruturada durante o desenvolvimento de jogos digitais. Adicionalmente, \citeonline{souza2023investigando} discutem dificuldades relacionadas à comunicação e à formalização do \textit{feedback} em equipes de desenvolvimento de \textit{software}. Nesse contexto, mecanismos de organização podem contribuir para reunir informações dispersas e facilitar sua consulta durante a tomada de decisão.

Considerando a dispersão das informações relacionadas ao desenvolvimento de jogos digitais, torna-se relevante investigar formas de reunir ideias, sugestões e \textit{feedbacks} em um ambiente organizado. Uma plataforma voltada a essa finalidade pode facilitar a consulta às contribuições dos jogadores e favorecer a interação entre a comunidade e os desenvolvedores.

Nesse contexto, a plataforma proposta busca apoiar a centralização e a organização de informações relacionadas a jogos digitais. Por meio de um protótipo funcional, o trabalho pretende demonstrar uma possibilidade de reunir essas contribuições de maneira mais acessível, auxiliando inicialmente a identificação de aspectos que possam ser considerados durante o desenvolvimento dos jogos.
